\clearpage
\section{Per-language tokenizer metrics}
\label{app:allthetokenizers}

\Cref{tab:intrinsic-quick} gives compression for every language and every scheme trained,
including the three that \Cref{tab:intrinsic-main} leaves out: case handling at the two
narrower boundary scopes, \bnd{wcaps} and \bnd{wpcaps}, and \bnd{wpdcapsin}, which places
the case code inside the marker pair rather than outside it. All three are controls, and
none of them changes the picture.

Case handling behaves the same way at every scope, which is why reporting it once is
enough. Adding it gains between 0.1 and 0.9 points of compression in English, German and
Russian under both trainers, is worth nothing in Korean ($-0.04$ at every scope), and
costs Arabic between 1.9 and 2.3 points. Finnish is the only language where the trainers
disagree about the sign: BPE loses 0.06 points and MinGram gains 0.9. The Arabic penalty
is flat across scopes and near-identical under both trainers, which makes it a fixed cost
of the mechanism rather than an interaction with the markers, but we cannot account for
it. It is not the cost of carrying case machinery through uncased text: Korean pays
nothing, and its evaluation slice holds proportionally more cased characters than Arabic's
(4.3\% against 1.5\%). Nor is it a vocabulary-budget effect: the case codes occupy two
slots out of 34{,}685, and the learned vocabulary is the same size either way.

Placement is a null result. The two placements differ by 0.03 points of compression in
favour of the inside one under BPE and by 0.18 in favour of the outside one under MinGram
-- the trainers disagree on the sign, and both differences are far smaller than the gap
between the trainers themselves. Neither morphology metric separates them either. What
does differ is vocabulary composition: in German under BPE, the number of learned
lowercase surfaces of at least three characters that also appear in an all-caps form falls
from 199 under \plainscheme{} to 146 under \bnd{wpd}, and to 62 and 56 under the inside
and outside placements respectively. Title-case duplication is left roughly where it was,
so what the case codes buy back is the all-caps variants specifically -- title-case forms
are frequent enough that the trainer keeps them regardless. We report the outside
placement because it leaves the marker pair a pure span delimiter, not because it
compresses better.

\begin{table}[h]
% Generated by marker_experiments/make_intrinsic_table.py. Do not edit.
% Requires booktabs and \newcommand{\bnd}[1]{\texttt{bnd\_#1}}.
% source: marker_experiments/paper/generated/eval_goldfish.json (quick corpus)
\centering
\small
\begin{tabular}{l rrrrr}
\toprule
& \multicolumn{5}{c}{BPE} \\
\cmidrule(lr){2-6}
Language & \texttt{plain} & \bnd{w} & \bnd{wp} & \bnd{wpd} & \bnd{wpd\_caps} \\
 & {\footnotesize ch/tok} & {\footnotesize \%} & {\footnotesize \%} & {\footnotesize \%} & {\footnotesize \%} \\
\midrule
German & 4.5734 & $-10.70$ & $-3.00$ & $-1.07$ & $-0.78$ \\
Arabic & 3.9677 & $-8.50$ & $-1.92$ & $-0.23$ & $-2.51$ \\
\midrule
Mean &  & $-9.60$ & $-2.46$ & $-0.65$ & $-1.65$ \\
\bottomrule
\end{tabular}
\caption{Compression, 5\,GB of FineWeb per language (the \texttt{quick} read-until-full sample, which is not uniform over the source), evaluated on held-out Goldfish, 34{,}685 matched total vocabulary. Baseline is \texttt{plain} in characters per token; each variant is the percentage change against that baseline. \bnd{w} $<$ \bnd{wp} $<$ \bnd{wpd} in every cell. The \texttt{plain} baseline itself fails to round-trip part of this slice (German 0.16\%, Arabic 3.32\% of documents), which is the base script encoding's handling of that text rather than an effect of the markers; no variant fails on more documents than its own baseline does.}
\label{tab:intrinsic-quick}

\end{table}
